\section{Summary and agent job descriptions}

Table \ref{tab:where-ai-fits} summarises the six steps.
Three patterns stand out.
Ordinary code has a part in every step, and running the model needs nothing else.
The LLM's part is reading, mapping, diagnosing and drafting, never computing.
And the person's part does not shrink to zero anywhere except the run: it moves from doing the work to reviewing it in the data steps, and stays a decision in the scenario, the calibration and the interpretation.

\begin{table}[htbp]
\centering
\small
\caption{Where code, an agent and a person fit in each step of GGM's workflow.}
\label{tab:where-ai-fits}
\begin{tabularx}{\textwidth}{@{}l>{\raggedright\arraybackslash}X>{\raggedright\arraybackslash}X>{\raggedright\arraybackslash}X@{}}
\toprule
Step & Ordinary code & Agent & Person \\
\midrule
Set up the scenario & Apply edits; compare workbooks & Find affected rows; propose edits and knock-on choices; check coherence & Question, storyline, node structure; approve edits \\
Get source data & Download and parse databases & Extract from reports; record source and confidence & Choose the anchoring outlook; review by confidence class \\
Prepare the data & Convert, aggregate, split, interpolate, check & Map sources to nodes; draft projections and demand parameters; explain warnings & Review mappings; decide projections and demand parameters \\
Calibrate & Deviation report; apply changes within bounds; run; log & Diagnose deviations; propose changes with rationale; research starting prices & Accept the calibration; choose between fits; market power \\
Run the model & Run the scenario matrix; record inputs and status & Explain failed runs & None routinely \\
Read the results & Compare scenarios; compute balances & Draft explanations; flag implausible patterns & Interpret and decide what to claim \\
\bottomrule
\end{tabularx}
\end{table}

Four agents follow from the table.
Each is described by what it is given, what it produces, the tools it uses, what it must never do, and how one would know that it works.
The common rule is that no number in a workbook or a result originates in the model's text: it comes from a source the agent cites, a function written and reviewed in advance, or a model run.

\textbf{A scenario agent.} It is given a storyline in plain language and the base workbooks, and produces a list of proposed edits, each with its reason, plus the knock-on choices a person should consider.
Its tools read and search the workbooks; code applies the approved edits and shows the difference from the base.
It must never change a workbook without approval or decide the storyline.
It works if, given the storylines behind the 2023 variants, it proposes the edits that distinguish them, without misses.

\textbf{A data agent.} It is given the node structure and a list of sources, and produces the general workbook and drafts of the projections and demand parameters, with the source and confidence class of every value.
Its tools retrieve documents, extract tables and call deterministic functions for conversion, aggregation and distances.
It must never do arithmetic in its own text or fill a gap without saying so.
It works if its general workbook matches one prepared by the developers, and every disagreement is either a newer source or a flagged uncertainty.

\textbf{A calibration agent.} It is given a deviation report, the parameters it may change with their bounds, and the documentation's diagnostic heuristics, and produces proposed changes with a written rationale for each.
Its tools are a harness that applies changes, runs the model locally and returns the next report.
It must never change elasticities or market power, exceed the bounds, or declare the calibration finished.
It works if it restores the fit of a deliberately perturbed 2023 scenario in less analyst time than calibration by hand, and its diagnoses agree with the documented heuristics.

\textbf{A results agent.} It is given the result tables, the comparisons computed by code and the scenario's inputs, and produces draft explanations and a list of patterns worth checking.
It must never state a number that is not in the tables.
It works if the analysts find that its flags point to real issues and its explanations hold when traced back to the tables.

Only the calibration agent is an agent in the full sense of Section 4.1, choosing its next action from what it observes; the other three are better described as workflows with an LLM at particular steps.
This answers the third research question: an agent could take over most of the reading and mapping and support the diagnosis in calibration, ordinary code should do all computation and running, and a person should remain responsible for the storyline, the projections, the acceptance of the calibration and the interpretation.
Chapter 6 discusses the risks of this division and which agent to build first.
